%\usepackage{tocbasic}      % already loaded by KOMA anyway
%\usepackage{titletoc}      % lighter alternative
%\usepackage{listings}
\usepackage{xcolor}
\usepackage[braket, qm]{qcircuit}
\usepackage{braket}
%\usepackage{etoc}                     % very powerful and fully compatible
\usepackage{microtype}
\usepackage{cancel}
\usepackage{subfigure}
\usepackage{enumitem}
\usepackage{tikz}
%\usepackage{quantikz}
%\usepackage{qcircuit}
%\usepackage{subcaption}
%\usepackage[utf8]{inputenc}
\usepackage{hyperref}            % لینک‌گذاری خودکار در PDF (ارجاع، فهرست، شکل‌ها)
\usepackage{tocloft}             % تنظیم فهرست مطالب و شکل‌ها
\usepackage{titlesec}            % کنترل ظاهر تیترها
%\usepackage[utf8]{inputenc}      % پشتیبانی از نویسه‌های UTF-8 (برای فارسی و لاتین)
\usepackage[T1]{fontenc}         % بهبود فونت و خروجی PDF
\usepackage{lmodern}             % فونت‌های بهتر لاتین
\usepackage{geometry}            % تنظیم حاشیه‌ها و اندازه صفحه
\usepackage{setspace}            % برای تنظیم فاصله خطوط (مثلاً 1.5)
\usepackage{graphicx}            % برای درج تصاویر
\usepackage{caption}             % تنظیم عنوان شکل‌ها و جدول‌ها
%\usepackage{subcaption}          % زیرشکل‌ها (subfigures)
\usepackage{booktabs}            % جدول‌های حرفه‌ای
\usepackage{array}               % ستون‌بندی بهتر در جدول‌ها
%\usepackage{mathtools}           % بهبود امکانات amsmath
\usepackage{physics}             % نمادهای رایج در فیزیک (bra-ket, derivatives, vectors,...)
%\usepackage{siunitx}             % نوشتن واحدها و مقادیر فیزیکی درست (SI)

% ░░░░░░░▒▒▒▒▒▒▓▓▓▓ Packages ▓▓▓▓▒▒▒▒▒▒░░░░░░░
\usepackage{amsthm,amssymb,mathtools}
\usepackage{dsfont}
\usepackage{amsfonts,fancybox,euscript,bbm,stmaryrd,eufrak}
% بسته‌‌ای برای چاپ شدن خودکار تعداد صفحات در صفحه «معرفی پایان‌نامه»
\usepackage{lastpage}
%\usepackage{amsthm}
\usepackage[nolist]{acronym}
% بسته‌‌ای برای ایجاد دیاگرام‌های مختلف
\usepackage[all]{xy}
%\usepackage{setspace}
%\usepackage{physics}
%%%%%%%%%%%%%%%%%%%%%%%%%%%%%%%%%%%%%%%%%%%%%%%%%%%%
% بسته‌‌ای برای ظاهر شدن شکل‌ها و تصاویر متن
%\usepackage{lastpage}
%\usepackage{adjustbox}
\usepackage{fancyhdr}
\usepackage{footnote}
\usepackage[numbers]{natbib}
%\usepackage{cite}
%\usepackage{tikz}
%\usepackage{tkz-graph}
% دستورات مربوط به ایجاد نمایه
%%%%%%%%%%%%%%%%%%%%%%%%%%%%%%%%%%%%%%%%%%%%%%%
\usepackage{zref-perpage}
%\usepackage{caption}% http://ctan.org/pkg/caption
%\captionsetup[figure]{labelformat=empty}%
\usepackage{pdfpages}
%تا اینجای پکیج‌ها از یه پایان‌نامه آورده شد...


\usepackage{upquote}
%\usepackage{listings}
%\usepackage{algorithm}
\usepackage[noend]{algpseudocode}
%\usepackage{unicode-math}
%\usepackage{amsmath,amssymb,amsthm,amsfonts}
%\usepackage{graphicx}
%\usepackage{geometry}
%\usepackage{color}
\usepackage{sectsty}
\usepackage{titlesec}
%\usepackage{caption}
%\usepackage[nottoc,notlof,notlot]{tocbibind}
%\usepackage{notoccite}
\usepackage[bottom]{footmisc}
%\usepackage{zref-perpage}
%\usepackage{tocloft}
\usepackage{array}
%\usepackage{enumitem}
%\usepackage{subfig}
%\usepackage{fancyhdr}
%\usepackage{indentfirst}
\usepackage{fontspec}
%\usepackage{setspace}
\usepackage{afterpage}
%\usepackage{fancybox}
%\usepackage[capitalise]{cleveref}% ارجاع هوشمند و خودکار به معادله، شکل و جدول
\usepackage{amsthm}
\theoremstyle{definition}
%\newtheorem{theorem}{قضیه}
\newtheorem{definition}{تعریف}
\newtheorem{lemma}{لم}
%\numberwithin{theorem}{section}
\newcommand{\esbat}{\medskip \noindent {\textbf{اثبات}:} }
\newcommand{\borhan}{\medskip \noindent {\textbf{برهان خلف}:} }


\newtheoremstyle{mystyle}%
{}%                                     
{}%                                    
{}%                             
{}%                                     
{\fontsize{12pt}{0}\selectfont \bfseries \BNazaninScaleOne} 
{ }%                                    
{ }%                                    
{ }%

\theoremstyle{mystyle}
\newtheorem{theorem}{قضیه}[section]
%\renewcommand\qedsymbol{$\blacksquare$}
%\newtheorem{remark}{تبصره}[chapter]
%\newtheorem{reminder}{یادآوری}[chapter]
\renewcommand{\thetheorem}{\thesection-\arabic{theorem}} % این خط شماره گذاری را به صورت دلخواه تنظیم میکند


\usepackage{xepersian}
\usepackage[free-standing-units]{siunitx}
\sisetup{per-mode = symbol}
% ░░░░░░░▒▒▒▒▒▒▓▓▓▓ Settings ▓▓▓▓▒▒▒▒▒▒░░░░░░░
\settextfont[Scale=1.2, ItalicFont=*, ItalicFeatures={FakeSlant=0.32}, BoldItalicFont=* Bold, BoldItalicFeatures={FakeSlant=0.32}]{IRNazli}
\setlatintextfont[Scale=1]{Times New Roman} % 11pt
\defpersianfont\BNazaninScaleOne[Scale=1]{IRNazli}
\defpersianfont\BTitrScaleOne[Scale=1]{IRNazli}
\deflatinfont\TimesNewRomanScaleOne[Scale=1]{Times New Roman}

\makeatletter
\def\abjad@zero{}
\def\abj@num@i#1{%
	\ifcase#1\or الف\or ب\or ج\or د%
	\or ه‍\or و\or ز\or ح\or ط\fi
	\ifnum#1=\z@\abjad@zero\fi}
\def\abj@num@ii#1{%
	\ifcase#1\or ی\or ک\or ل\or م\or ن%
	\or س\or ع\or ف\or ص\fi
	\ifnum#1=\z@\fi\abj@num@i}
\def\abj@num@iii#1{%
	\ifcase#1\or ق\or ر\or ش\or ت\or ث%
	\or خ\or ذ\or ض\or ظ\fi
	\ifnum#1=\z@\fi\abj@num@ii}
\def\abj@num@iv#1{%
	\ifcase#1\or غ\fi
	\ifnum#1=\z@\fi\abj@num@iii}
\let\@latinalph\@alph%
\let\@latinAlph\@Alph%
\def\PersianAlphs{%
	\let\@alph\abjadnumeral%
	\let\@Alph\abjadnumeral}
\def\LatinAlphs{%
	\let\@alph\@latinalph%
	\let\@Alph\@latinAlph}
\PersianAlphs
\makeatother

%%
\setlength\intextsep{2\baselineskip plus 2mm minus 1mm}
\setlength\textfloatsep{3\baselineskip plus 2mm minus 1mm}
\setlength\floatsep{3\baselineskip}
\setlength{\tabcolsep}{1.5mm}
\renewcommand{\arraystretch}{1.5}
\DeclareCaptionFont{MyFont}{\fontsize{12pt}{0}\selectfont \bfseries \BNazaninScaleOne} % B Nazanin 12
\DeclareCaptionLabelSeparator{MySeparator}{\,-\,}
\captionsetup{margin=5mm, font=MyFont, labelsep=MySeparator}

\captionsetup[figure]{aboveskip=2mm}
\captionsetup[figure]{belowskip=-\baselineskip}
\captionsetup[table]{aboveskip=2mm}
\captionsetup[table]{belowskip=\baselineskip}

%%
%\setmathfont[version=lm]{Latin Modern Math}
%\mathversion{lm}
\DeclareMathSizes{10}{11}{10}{8} % Font 12
%\DefaultMathsDigits
%\setlength{\mathindent}{0.5cm}

\AtBeginDocument{
	\abovedisplayskip=\baselineskip
	\belowdisplayskip=\baselineskip
	\abovedisplayshortskip=\baselineskip
	\belowdisplayshortskip=\baselineskip
}

\setlength\mathsurround{1pt}
\delimitershortfall=-3pt

%%
\geometry{top=3.5cm, bottom=3.5cm, left=3.5cm, right=3.5cm, headsep=1cm}

%%
\definecolor{darkblue}{rgb}{0,0,0.5}
\definecolor{darkred}{rgb}{0.5,0,0}
\hypersetup{colorlinks=true, linkcolor=darkblue, citecolor=darkred, urlcolor=darkblue, bookmarks, bookmarksdepth=3}

%%
\SepMark{-}

%%
%\renewcommand*{\citedash}{-}

\makeatletter
\def\@cite#1#2{[{#1\if@tempswa ,~#2\fi}]}
\makeatother

%%
\setlength{\parindent}{5mm}
%
%%%
%\newtheoremstyle{mystyle}%
%{}%                                     
%{}%                                    
%{\itshape}%                             
%{}%                                     
%{\fontsize{12pt}{0}\selectfont \bfseries \BNazaninScaleOne} 
%{-}%                                    
%{ }%                                    
%{}%
%
%\theoremstyle{mystyle}
%\newtheorem{theorem}{قضیه}[subsection]
%%\renewcommand\qedsymbol{$\blacksquare$}
%%\newtheorem{remark}{تبصره}[chapter]
%%\newtheorem{reminder}{یادآوری}[chapter]

%%
\titleformat{\chapter}[display]
{\fontsize{15pt}{0}\selectfont  \flushleft  \bfseries \BNazaninScaleOne} % B Nazanin 15
{\chaptertitlename\ \tartibi{chapter}}{0.3cm}{\fontsize{15pt}{0}\selectfont \BNazaninScaleOne} % B Nazanin 15

\sectionfont{\fontsize{14pt}{\baselineskip}\selectfont \bfseries \BNazaninScaleOne} % B Nazanin 14
\subsectionfont{\fontsize{13pt}{\baselineskip}\selectfont \bfseries \BNazaninScaleOne} % B Nazanin 13
\subsubsectionfont{\fontsize{13pt}{\baselineskip}\selectfont \bfseries \BNazaninScaleOne} % B Nazanin 13

\titlespacing*{\chapter}{0pt}{3.5cm}{8cm}
\titlespacing*{\section}{0pt}{1.2\baselineskip}{0pt}
\titlespacing*{\subsection}{0pt}{1.1\baselineskip}{0pt}
\titlespacing*{\subsubsection}{0pt}{1\baselineskip}{0pt}

\makeatletter
\renewcommand\chapter{\if@openright\cleardoublepage\else\clearpage\fi
	\thispagestyle{empty}%
	\global\@topnum\z@
	\@afterindentfalse
	\secdef\@chapter\@schapter}
\makeatother

\setcounter{secnumdepth}{3}

%%
\setlist[enumerate]{noitemsep, topsep=0pt}
\setlist[itemize]{noitemsep, topsep=0pt}
\setlist[description]{noitemsep, topsep=0pt}

%%
\zmakeperpage{footnote}
\setlength{\footnotesep}{0.4cm}
\renewcommand{\footnotesize}{\fontsize{8.5pt}{0}\selectfont} 
\renewcommand{\thefootnote}{\fontsize{6.5pt}{0}\selectfont\arabic{footnote}}
\setlength{\footnotemargin}{0.2cm}

%%
\setcounter{tocdepth}{2}
\setlength{\cftbeforetoctitleskip}{0cm}
\setlength{\cftbeforeloftitleskip}{0cm}
\setlength{\cftbeforelottitleskip}{0cm}

\setlength{\cftaftertoctitleskip}{1cm}
\setlength{\cftafterloftitleskip}{1cm}
\setlength{\cftafterlottitleskip}{1cm}

\renewcommand{\cfttoctitlefont}{\hspace{6.5cm} \fontsize{13pt}{0}\selectfont \bfseries \BNazaninScaleOne} % B Nazanin 13
\renewcommand{\cftloftitlefont}{\hspace{6.5cm} \fontsize{13pt}{0}\selectfont \bfseries \BNazaninScaleOne} % B Nazanin 13
\renewcommand{\cftlottitlefont}{\hspace{6.5cm} \fontsize{13pt}{0}\selectfont \bfseries \BNazaninScaleOne} % B Nazanin 13

\renewcommand{\listfigurename}{\textbf{\fontsize{13pt}{0}\selectfont \BNazaninScaleOne} فهرست شکل‌ها}
\renewcommand{\listtablename}{\textbf{\fontsize{13pt}{0}\selectfont \BNazaninScaleOne} فهرست جدول‌ها}

\renewcommand{\cftfigpresnum}{شکل }
\renewcommand{\cftfigaftersnum}{:}

\renewcommand{\cfttabpresnum}{جدول }
\renewcommand{\cfttabaftersnum}{:}

\renewcommand{\cftchapfont}{\bfseries \fontsize{13pt}{0}\selectfont \BNazaninScaleOne} % B Nazanin 13
\renewcommand{\cftsecfont}{\fontsize{13pt}{0}\selectfont \BNazaninScaleOne} % B Nazanin 13
\renewcommand{\cftsubsecfont}{\fontsize{13pt}{0}\selectfont \BNazaninScaleOne} % B Nazanin 13
\renewcommand{\cftsubsubsecfont}{\fontsize{13pt}{0}\selectfont \BNazaninScaleOne} % B Nazanin 13
\renewcommand{\cftfigfont}{\fontsize{13pt}{0}\selectfont \BNazaninScaleOne} % B Nazanin 13
\renewcommand{\cfttabfont}{\fontsize{13pt}{0}\selectfont \BNazaninScaleOne} % B Nazanin 13

\renewcommand\cftchappagefont{\fontsize{13pt}{0}\selectfont \BNazaninScaleOne} % B Nazanin 13
\renewcommand\cftsecpagefont{\fontsize{13pt}{0}\selectfont \BNazaninScaleOne} % B Nazanin 13
\renewcommand\cftsubsecpagefont{\fontsize{13pt}{0}\selectfont \BNazaninScaleOne} % B Nazanin 13
\renewcommand\cftsubsubsecpagefont{\fontsize{13pt}{0}\selectfont \BNazaninScaleOne} % B Nazanin 13
\renewcommand\cftfigpagefont{\fontsize{13pt}{0}\selectfont \BNazaninScaleOne} % B Nazanin 13
\renewcommand\cfttabpagefont{\fontsize{13pt}{0}\selectfont \BNazaninScaleOne} % B Nazanin 13

\renewcommand\cftchapafterpnum{\vspace{0cm}}
\renewcommand\cftsecafterpnum{\vspace{0cm}}
\renewcommand\cftsubsecafterpnum{\vspace{0cm}}
\renewcommand\cftsubsubsecafterpnum{\vspace{0cm}}

\renewcommand{\cftdotsep}{0.2}

\setlength\cftchapindent{0cm}
\setlength\cftsecindent{0cm}
\setlength\cftsubsecindent{0.35cm}
\setlength\cftsubsubsecindent{0cm}
\setlength\cftfigindent{0cm}
\setlength\cfttabindent{0cm}

\makeatletter
\renewcommand{\numberline}[1]{%
	\@cftbsnum #1\@cftasnum \hspace*{2mm}\@cftasnumb%
}
\makeatother

\addtocontents{toc}{~\textbf{\fontsize{13pt}{0}\selectfont {\BNazaninScaleOne عنوان}}\hfill\textbf{\fontsize{13pt}{0}\selectfont {\BNazaninScaleOne صفحه}}\vspace{0.2cm}\protect\par} % B Nazanin 13
\addtocontents{toc}{\protect\vspace*{15pt}}
%\addtocontents{toc}{\protect\afterpage{\hspace*{-3.9em}\bfseries عنوان \hfill صفحه}}
\addtocontents{toc}{\protect\afterpage{\hspace*{-1.3em}\bfseries عنوان \hfill صفحه}}
\addtocontents{toc}{\protect\afterpage{\vspace*{25pt}}}

\addtocontents{lof}{~\textbf{\fontsize{13pt}{0}\selectfont {\BNazaninScaleOne عنوان}}\hfill\textbf{\fontsize{13pt}{0}\selectfont {\BNazaninScaleOne صفحه}}\vspace{-0.1cm}\protect\par} % B Nazanin 13
\addtocontents{lof}{\protect\vspace*{15pt}}
\addtocontents{lof}{\protect\afterpage{\hspace*{-1.3em}\bfseries عنوان \hfill صفحه}}
\addtocontents{lof}{\protect\afterpage{\vspace*{25pt}}}

\addtocontents{lot}{~\textbf{\fontsize{13pt}{0}\selectfont {\BNazaninScaleOne عنوان}}\hfill\textbf{\fontsize{13pt}{0}\selectfont {\BNazaninScaleOne صفحه}}\vspace{-0.1cm}\protect\par} % B Nazanin 13
\addtocontents{lot}{\protect\vspace*{15pt}}
\addtocontents{lot}{\protect\afterpage{\hspace*{-1.3em}\bfseries عنوان \hfill صفحه}}
\addtocontents{lot}{\protect\afterpage{\vspace*{25pt}}}

\renewcommand\cftchapaftersnum{:}
\makeatletter
\bidi@patchcmd{\Hy@org@chapter}{%
	\addcontentsline{toc}{chapter}%
	{\protect\numberline{\thechapter}#1}%
}{%
	\addcontentsline{toc}{chapter}%
	{\protect\numberline{\chaptername~\tartibi{chapter}}#1}%
}{}{}

\makeatother

%%
\makeatletter

\def\DepartmentFa#1{\gdef\@DepartmentFa{#1}}
\def\@DepartmentFa{}

\def\GroupFa#1{\gdef\@GroupFa{#1}}
\def\@GroupFa{}

\def\TitleFa#1{\gdef\@TitleFa{#1}}
\def\@TitleFa{}

\def\ThesisTypeFa#1{\gdef\@ThesisTypeFa{#1}}
\def\@ThesisTypeFa{}

\def\DegreeFa#1{\gdef\@DegreeFa{#1}}
\def\@DegreeFa{}

\def\ThesisMark#1{\gdef\@ThesisMark{#1}}
\def\@ThesisMark{}


\def\FieldFa#1{\gdef\@FieldFa{#1}}
\def\@FieldFa{}

\def\BranchFa#1{\gdef\@BranchFa{#1}}
\def\@BranchFa{}

\def\YourFullnameFa#1{\gdef\@YourFullnameFa{#1}}
\def\@YourFullnameFa{}

\def\YourEmailAddress#1{\gdef\@YourEmailAddress{#1}}
\def\@YourEmailAddress{}

\def\FirstSupervisorFa#1{\gdef\@FirstSupervisorFa{#1}}
\def\@FirstSupervisorFa{}
\def\FirstSupervisorAcademicRank#1{\gdef\@FirstSupervisorAcademicRank{#1}}
\def\@FirstSupervisorAcademicRank{}

\def\SecondSupervisorFa#1{\gdef\ST@SecondSupervisorFa{#1}}
\def\@SecondSupervisorFa{}
\def\SecondSupervisorAcademicRank#1{\gdef\@SecondSupervisorAcademicRank{#1}}
\def\@SecondSupervisorAcademicRank{}

\def\FirstAdvisorFa#1{\gdef\@FirstAdvisorFa{#1}}
\def\@FirstAdvisorFa{}
\def\FirstAdvisorAcademicRank#1{\gdef\@FirstAdvisorAcademicRank{#1}}
\def\@FirstAdvisorAcademicRank{}

\def\SecondAdvisorFa#1{\gdef\ST@SecondAdvisorFa{#1}}
\def\@SecondAdvisorFa{}
\def\SecondAdvisorAcademicRank#1{\gdef\@SecondAdvisorAcademicRank{#1}}
\def\@SecondAdvisorAcademicRank{}

\def\FirstExaminerFa#1{\gdef\ST@FirstExaminerFa{#1}}
\def\@FirstExaminerFa{}
\def\FirstExaminerAcademicRank#1{\gdef\@FirstExaminerAcademicRank{#1}}
\def\@FirstExaminerAcademicRank{}

\def\SecondExaminerFa#1{\gdef\ST@SecondExaminerFa{#1}}
\def\@SecondExaminerFa{}
\def\SecondExaminerAcademicRank#1{\gdef\@SecondExaminerAcademicRank{#1}}
\def\@SecondExaminerAcademicRank{}

\def\ThirdExaminerFa#1{\gdef\ST@ThirdExaminerFa{#1}}
\def\@ThirdExaminerFa{}
\def\ThirdExaminerAcademicRank#1{\gdef\@ThirdExaminerAcademicRank{#1}}
\def\@ThirdExaminerAcademicRank{}

\def\DeanOfDepartmentFa#1{\gdef\@DeanOfDepartmentFa{#1}}
\def\@DeanOfDepartmentFa{}

\def\YearFa#1{\gdef\@YearFa{#1}}
\def\@YearFa{}

\def\DateFa#1{\gdef\@DateFa{#1}}
\def\@DateFa{}

\def\Prefix#1{\gdef\@Prefix{#1}}
\def\@Prefix{}

\def\DepartmentEn#1{\gdef\@DepartmentEn{#1}}
\def\@DepartmentEn{}

\def\GroupEn#1{\gdef\@GroupEn{#1}}
\def\@GroupEn{}

\def\TitleEn#1{\gdef\@TitleEn{#1}}
\def\@TitleEn{}

\def\DegreeEn#1{\gdef\@DegreeEn{#1}}
\def\@DegreeEn{}

\def\YourFullnameEn#1{\gdef\@YourFullnameEn{#1}}
\def\@YourFullnameEn{}

\def\DateEn#1{\gdef\@DateEn{#1}}
\def\@DateEn{}

\def\FirstSupervisorEn#1{\gdef\@FirstSupervisorEn{#1}}
\def\@FirstSupervisorEn{}

\def\SecondSupervisorEn#1{\gdef\ST@SecondSupervisorEn{#1}}
\def\@SecondSupervisorEn{}

\def\FirstAdvisorEn#1{\gdef\@FirstAdvisorEn{#1}}
\def\@FirstAdvisorEn{}

\def\SecondAdvisorEn#1{\gdef\ST@SecondAdvisorEn{#1}}
\def\@SecondAdvisorEn{}

\def\FirstExaminerEn#1{\gdef\ST@FirstExaminerEn{#1}}
\def\@FirstExaminerEn{}

\def\SecondExaminerEn#1{\gdef\ST@SecondExaminerEn{#1}}
\def\@SecondExaminerEn{}

\def\ThirdExaminerEn#1{\gdef\ST@ThirdExaminerEn{#1}}
\def\@ThirdExaminerEn{}

\def\FourthExaminerEn#1{\gdef\ST@FourthExaminerEn{#1}}
\def\@FourthExaminerEn{}

\def\FifthExaminerEn#1{\gdef\ST@FifthExaminerEn{#1}}
\def\@FifthExaminerEn{}

\def\DeanOfDepartmentEn#1{\gdef\@DeanOfDepartmentEn{#1}}
\def\@DeanOfDepartmentEn{}

\def\FirstSupervisorEmailAddress#1{\gdef\@FirstSupervisorEmailAddress{#1}}
\def\@FirstSupervisorEmailAddress{}

\def\SecondSupervisorEmailAddress#1{\gdef\@SecondSupervisorEmailAddress{#1}}
\def\@SecondSupervisorEmailAddress{}

\makeatother

% ░░░░░░░▒▒▒▒▒▒▓▓▓▓ Title Page ▓▓▓▓▒▒▒▒▒▒░░░░░░░
\makeatletter
\newcommand{\MakeTitlePage}{
	\clearpage
	\thispagestyle{empty}
	\newgeometry{left=3.5cm, right=3.5cm, top=3.5cm, bottom=3.5cm}
	{\BNazaninScaleOne
		\begin{center}
			\includegraphics[width = 2.5cm]{UI_Logo.pdf}
			\\[0.5cm]
			{\fontsize{13pt}{0}\selectfont % B Nazanin bold 13
				\textbf{
					دانشگاه اصفهان
			}}\\[0.35cm]
			{\fontsize{12pt}{0}\selectfont \textbf{\@DepartmentFa}}\\ [0.3cm] % B Nazanin bold 12
			{\fontsize{11pt}{0}\selectfont \textbf{\@GroupFa}}\\ % B Nazanin bold 11
			\vspace*{\fill}
			{\fontsize{18pt}{0}\selectfont \textbf{\@ThesisTypeFa~\@DegreeFa~\@FieldFa}}\\[1.5cm] % B Nazanin bold 18
			{\fontsize{15pt}{0}\selectfont  \BTitrScaleOne \@TitleFa}\\ % B Titr 15
			\vspace*{\fill}
			
			\ifx\ST@SecondSupervisorFa\undefined
			{\fontsize{15pt}{0}\selectfont \textbf{
					استاد راهنما:
			}} \else {\fontsize{15pt}{0}\selectfont \textbf{
					استادهای راهنما:
			}}\fi \\[0.4cm] % B Nazanin 15
			\ifx\ST@SecondSupervisorFa\undefined
			{\fontsize{13pt}{0}\selectfont \textbf{دکتر‍~\@FirstSupervisorFa}}
			\else {\fontsize{13pt}{0}\selectfont \textbf{\@FirstSupervisorFa~~و~~\ST@SecondSupervisorFa}} \fi \\% B Nazanin  bold 13
			\vspace*{\fill}
			
			{\fontsize{15pt}{0}\selectfont \textbf{
					دانشجو:
			}}  \\[0.4cm] % B Nazanin bold 15
			
			{\fontsize{13pt}{0}\selectfont \textbf{\@YourFullnameFa}}\\ % ‌B Nazanin bold 13
			\vspace*{\fill}
			
			{\fontsize{13pt}{0pt}\selectfont \textbf{\@YearFa}} % B Nazanin bold 13
		\end{center}
	}
	\restoregeometry
}

\makeatother

% ░░░░░░░▒▒▒▒▒▒▓▓▓▓ CopyRight Page ▓▓▓▓▒▒▒▒▒▒░░░░░░░
\newcommand{\MakeCopyRightPage}{
	\clearpage
	\thispagestyle{empty}
	\vspace*{\fill}
	\hspace*{\fill}
	\begin{minipage}{11cm}
		{\BNazaninScaleOne
			{\fontsize{13pt}{0.9cm}\selectfont  %‌ B Nazanin 13
				\noindent
				{\bfseries
					کليه‌ی حقوق مادی و معنوی مترتب بر دستاوردهای مطالعات، ابتکارات و نوآوری‌های ناشی از پژوهش موضوع اين رساله متعلق به دانشگاه اصفهان است. دانشجو موظف به رعایت آیین‌نامه و منشور اخلاق در پژوهش برای ارائه و یا چاپ مطالب مستخرج از رساله خود می‌باشد.
					
		}}}
	\end{minipage}
	\hspace*{\fill}
	\vspace*{\fill}
}

% ░░░░░░░▒▒▒▒▒▒▓▓▓▓ Signature Page - Farsi ▓▓▓▓▒▒▒▒▒▒░░░░░░░
\newcounter{RowNoFa}
\newcommand\RowNumberFa{\stepcounter{RowNoFa}\arabic{RowNoFa}}

\makeatletter
\newcommand{\MakeFarsiSignaturePage}{
	\clearpage
	\thispagestyle{empty}
	\newgeometry{left=3.5cm, right=3.5cm, top=3.5cm, bottom=3.5cm}
	{\BNazaninScaleOne
		\begin{center}
			\includegraphics[width = 2.5cm]{UI_Logo.pdf}
			\\ [0.5cm]
			{\fontsize{13pt}{0}\selectfont % B Nazanin bold 13
				\textbf{
					دانشگاه اصفهان
			}}\\[0.35cm]
			{\fontsize{12pt}{0}\selectfont \textbf{\@DepartmentFa}}\\ [0.3cm] % B Nazanin bold 12
			{\fontsize{11pt}{0}\selectfont \textbf{\@GroupFa}}\\ % B Nazanin bold 11
			\vspace{3cm}
			
			{\fontsize{18pt}{0}\selectfont \textbf{\@ThesisTypeFa~\@DegreeFa~رشته‌ی~\@FieldFa~گرایش~\@BranchFa~ \\ [0.4cm] \@Prefix~\@YourFullnameFa~تحت عنوان}}\\ % B Nazanin bold 18
			[1cm]
			
			{\fontsize{15pt}{0}\selectfont  \BTitrScaleOne \textbf{\@TitleFa}} % B Titr 15
			\\ [1cm]
			
		\end{center}
		{\fontsize{13pt}{0}\selectfont % B Nazanin 13
			در تاریخ
			~\@DateFa~
			توسط هیأت داوران بررسی و با درجه‌ی~\textbf{\@ThesisMark}~به تصویب نهایی رسید.‎\\
			\vfill
			
			\noindent
			\begin{tabular}{rrrrrr}
				\RowNumberFa - &
				استاد راهنمای~\@ThesisTypeFa & دکتر~\@FirstSupervisorFa & با مرتبه‌ی علمی~\textbf{\@FirstSupervisorAcademicRank} & امضاء
				\\[1cm]
				\ifx\ST@SecondSupervisorFa\undefined \else
				\RowNumberFa - &
				استاد راهنمای~\@ThesisTypeFa & دکتر~\@SecondSupervisorFa & با مرتبه‌ی علمی~\textbf{\@SecondSupervisorAcademicRank} &   امضاء
				%\\[1cm]
				\fi
				%\RowNumberFa - &
				%استاد مشاور~\@ThesisTypeFa  & دکتر‍~\ST@FirstSupervisorFa & با مرتبه‌ی علمی~\textbf{\@FirstSupervisorAcademicRank} &   امضاء
				%\\[1cm]
				\ifx\ST@SecondAdvisorFa\undefined \else
				%\RowNumberFa - &
				%استاد مشاور~\@ThesisTypeFa  & دکتر‍~\ST@SecondSupervisorFa & با مرتبه‌ی علمی~\textbf{\@SecondSupervisorAcademicRank}   & امضاء
				%\\[1cm]
				\fi
				\ifx\ST@FirstExaminerFa\undefined \else
				\RowNumberFa - &
				استاد داور اول داخل گروه  & دکتر‍~\ST@FirstExaminerFa & با مرتبه‌ی علمی~\textbf{\@FirstExaminerAcademicRank} & امضاء
				\\[1cm]
				\fi
				\ifx\ST@SecondExaminerFa\undefined \else
				\RowNumberFa - &
				استاد داور دوم داخل گروه  & دکتر‍~\ST@SecondExaminerFa & با مرتبه‌ی علمی~\textbf{\@SecondExaminerAcademicRank} & امضاء
				\\[1cm]
				\fi
				\ifx\ST@ThirdExaminerFa\undefined \else
				\RowNumberFa - &
				استاد داور خارج از گروه  & دکتر‍~\ST@ThirdExaminerFa & با مرتبه‌ی علمی~\textbf{\@ThirdExaminerAcademicRank} & امضاء
				\\[1cm]
				\fi
				
			\end{tabular}
		}
	}
	
	\vspace{1.5cm}
	\hfill{مهر و امضاء مدیرگروه}
	\restoregeometry
}

\makeatother

% ░░░░░░░▒▒▒▒▒▒▓▓▓▓ Farsi Abstract Page ▓▓▓▓▒▒▒▒▒▒░░░░░░░
%\makeatletter
%\def\AbstractFa#1{\gdef\@AbstractFa{#1}}
%\def\@AbstractFa{}
%
%\def\KeywordsFa#1{\gdef\@KeywordsFa{#1}}
%\def\@KeywordsFa{}
%
%\newcommand{\MakeFarsiAbstract}{
	%\clearpage
	%\thispagestyle{empty}
	%%\phantomsection
	%%\vspace*{\fill}
	%%\addcontentsline{toc}{section}{چکیده}‎
	%\baselineskip=0.7cm
	%{\fontsize{13pt}{0}\selectfont \BNazaninScaleOne \par \noindent \textbf{چکیده}} \\
	%[0.4cm]
	%{\fontsize{13pt}{0}\selectfont \BNazaninScaleOne \par \noindent \@AbstractFa} \\
	%[0.7cm]
	%{\fontsize{13pt}{0}\selectfont \BNazaninScaleOne \par \noindent \textbf{کلیدواژه‌ها:}} \\ % B Nazanin 13
	%[0.4cm]
	%{\fontsize{12pt}{0}\selectfont \BNazaninScaleOne \par\noindent \@KeywordsFa} % B Nazanin 12
	%}
%\makeatother

\makeatletter
\def\AbstractFa#1{\gdef\@AbstractFa{#1}}
\def\@AbstractFa{}

\def\KeywordsFa#1{\gdef\@KeywordsFa{#1}}
\def\@KeywordsFa{}
\makeatother

\makeatletter
\newcommand{\MakeFarsiAbstract}{
	\clearpage
	\thispagestyle{empty}
	\phantomsection
	%\vspace*{\fill}
	%\addcontentsline{toc}{section}{چکیده}‎
	\baselineskip=0.7cm
	{\BNazaninScaleOne
		{\fontsize{13pt}{0}\selectfont \par\noindent \textbf{چکیده} \baselineskip=0.7cm \\ \@AbstractFa \\ \textbf{کلمات کلیدی:} \baselineskip=0.7cm \\ \@KeywordsFa} % B Nazanin 13
	}
}
\makeatother

% ░░░░░░░▒▒▒▒▒▒▓▓▓▓ Table of Contents ▓▓▓▓▒▒▒▒▒▒░░░░░░░
\newcommand{\MakeTableOfContents}{
	\clearpage
	\newgeometry{top=3.5cm, left=3.5cm, right=3.5cm, bottom=3.5cm, footskip=1cm}
	\baselineskip=0.6cm
	\phantomsection
	\addcontentsline{toc}{section}{فهرست مطالب}‎
	\tableofcontents
	\restoregeometry
}

% ░░░░░░░▒▒▒▒▒▒▓▓▓▓ List of Figures ▓▓▓▓▒▒▒▒▒▒░░░░░░░
\newcommand{\MakeListOfFigures}{
	\clearpage
	\newgeometry{top=3.5cm, left=3.5cm, right=3.5cm, bottom=3.5cm, footskip=1cm}
	\baselineskip=0.6cm
	\phantomsection
	\addcontentsline{toc}{section}{\listfigurename}‎
	\listoffigures
	\restoregeometry
}

% ░░░░░░░▒▒▒▒▒▒▓▓▓▓ List of Tables ▓▓▓▓▒▒▒▒▒▒░░░░░░░
\newcommand{\MakeListOfTables}{
	\clearpage
	\newgeometry{top=3.5cm, left=3.5cm, right=3.5cm, bottom=3.5cm, footskip=1cm}
	\baselineskip=0.6cm
	\phantomsection
	\addcontentsline{toc}{section}{\listtablename}‎
	\listoftables
	\restoregeometry
}

\makeatother

% ░░░░░░░▒▒▒▒▒▒▓▓▓▓ English Abstract Page ▓▓▓▓▒▒▒▒▒▒░░░░░░░
\makeatletter
\def\AbstractEn#1{\gdef\@AbstractEn{#1}}
\def\@AbstractEn{}

\def\KeywordsEn#1{\gdef\@KeywordsEn{#1}}
\def\@KeywordsEn{}

\newcommand{\MakeEnglishAbstract}{
	\clearpage
	\thispagestyle{empty}
	\begin{latin}
		\baselineskip=0.6cm
		{\fontsize{12pt}{0}\selectfont \par\noindent \textbf{Abstract}\\ [0.4cm]\@AbstractEn \\ [0.4cm] \textbf{Keywords:}\\ [0.4cm]\@KeywordsEn} % TNR 12
	\end{latin}
}
\makeatother

% ░░░░░░░▒▒▒▒▒▒▓▓▓▓ Signature Page - English ▓▓▓▓▒▒▒▒▒▒░░░░░░░
\newcounter{RowNoEn}
\newcommand\RowNumberEn{\stepcounter{RowNoEn}\arabic{RowNoEn}}

\makeatletter
\newcommand{\MakeEnglishSignaturePage}{
	\clearpage
	\thispagestyle{empty}
	\newgeometry{left=3.5cm, right=3.5cm, top=3.5cm, bottom=3.5cm}
	\begin{latin}
		{\TimesNewRomanScaleOne
			\begin{center}
				\includegraphics[width=2.5cm]{UI_Logo.pdf}
				\\ [0.5cm]
				{\fontsize{12pt}{0}\selectfont \textbf{University of Isfahan}}\\[0.2cm] % TNR 12
				{\fontsize{11pt}{0}\selectfont \textbf{{\@DepartmentEn}}}\\ [0.2cm] % TNR 11
				{\fontsize{10pt}{0}\selectfont \textbf{{\@GroupEn}}}\\ % TNR 10
				\vspace*{\fill}
				{\fontsize{16pt}{0}\selectfont \textbf{\@DegreeEn~Thesis}}\\ % TNR 16
				\vspace*{\fill}
				
				{\fontsize{15pt}{0}\selectfont \textbf{\@TitleEn}} % TNR 15
				\vspace*{\fill}
				
				{\fontsize{14pt}{0}\selectfont \textbf{Supervisor:}} % TNR 14
				{\fontsize{13pt}{0}\selectfont \textbf{\\[0.25cm] Dr.~\@FirstSupervisorEn}} % TNR 13
				\vspace*{\fill}
				
				{\fontsize{14pt}{0}\selectfont \textbf{By:\\[0.25cm]}} % TNR 14
				{\fontsize{13pt}{0}\selectfont \textbf{\@YourFullnameEn}} % TNR 13
				\vspace*{\fill}
				
				{\fontsize{13pt}{0}\selectfont \textbf{\@DateEn}} % TNR 13
				\vspace*{\fill}
				
			\end{center}
		}
	\end{latin}
	\restoregeometry
}
\makeatother

% ░░░░░░░▒▒▒▒▒▒▓▓▓▓ Appendices ▓▓▓▓▒▒▒▒▒▒░░░░░░░
\newcommand{\MakeAppendices}{
	\clearpage
	
	\appendix
	
	\setcounter{chapter}{3}
	\setcounter{table}{0}
	\setcounter{figure}{0}
	\setcounter{equation}{0}
	\setcounter{theorem}{0}
	\setcounter{remark}{0}
	\setcounter{reminder}{0}
	
	\phantomsection
	
	\addtocontents{toc}{\protect\setcounter{tocdepth}{-1}}
	\chapter*{پیوست‌ها}
	\addtocontents{toc}{\protect\setcounter{tocdepth}{2}}
	
	\addcontentsline{toc}{chapter}{پیوست‌ها}‎
}

\makeatother
% ░░░░░░░▒▒▒▒▒▒▓▓▓▓ References ▓▓▓▓▒▒▒▒▒▒░░░░░░░
\newcommand{\MakeReferences}{
	\clearpage
	%\phantomsection
	%\section*{مراجع}
	%\addcontentsline{toc}{chapter}{مراجع}
	\titlespacing*{\chapter}{0pt}{-\baselineskip}{0.5cm}
	\fontsize{11pt}{14.4}\selectfont \BNazaninScaleOne \renewcommand{\bibname}{{مراجع}\hfill}
}

\graphicspath{{Figures/}}

%
%\theoremstyle{definition}
%\newtheorem{definition}{تعریف}[section]
%%\theoremstyle{theorem}
%\newtheorem{theorem}[definition]{قضیه}
%\newtheorem{lemma}[definition]{لم}
%\newtheorem{proposition}[definition]{گزاره}
%\newtheorem{corollary}[definition]{نتیجه}
%\newtheorem{remark}[definition]{ملاحظه}
%\theoremstyle{definition}
%\newtheorem{example}[definition]{مثال}
%\newtheorem{nok}[definition]{نکته}
%\newtheorem{question}[definition]{سوال}
%\newcommand{\esbat}{\medskip \noindent {\it اثبات:} }
%\newtheorem{theorem}{قضیه}
%\usepackage{amsthm}
%\newtheorem{theorem}{قضیه}[section]
%\newtheorem{definition}{تعریف}[section]



%
%\usepackage{xcolor}
%\usepackage{listings}
%
%\lstset{
%	basicstyle=\ttfamily\footnotesize,
%%	backgroundcolor=\color{black!5},
%	frame=single,
%	breaklines=true,
%	numbers=left,
%	numberstyle=\tiny\color{gray},
%	xleftmargin=10pt,
%	literate={-}{{-}}1,
%}
\lstset{
	basicstyle=\ttfamily\small, % monospaced font
	backgroundcolor=\color{black!3},   % light gray background
	frame=none,                        % <-- remove border lines
	numbers=none,                      % <-- no line numbers
	showstringspaces=false,
	breaklines=true,
}

%\lstset{
%	frame=none,          % بدون قاب دور کد
%	numbers=none,        % بدون شماره‌گذاری خطوط
%	xleftmargin=0pt,     % حذف حاشیه چپ
%	breaklines=true,     % شکستن خودکار خطوط طولانی
%	breakindent=0pt,     % عدم تورفتگی برای خطوط شکسته
%	postbreak=\mbox{\textcolor{gray}{$\hookrightarrow$}\space}, % نشانه برای خطوط شکسته شده
%	backgroundcolor=\color{white}, % رنگ پس‌زمینه کد (اختیاری)
%	keywordstyle=\bfseries, % استایل کلمات کلیدی
%	commentstyle=\itshape\color{gray}, % استایل کامنت‌ها
%	showstringspaces=false, % عدم نمایش فاصله در رشته‌ها	basicstyle=\ttfamily\small % استایل فونت اصلی کد
%}